\documentclass{article}
\usepackage{amsmath}
\usepackage{amssymb}
\usepackage[margin=1in]{geometry}

\newcommand{\sixj}[6]{\left\{\begin{matrix}#1&#2&#3\\#4&#5&#6\end{matrix}\right\}}

\title{Implemented Equations From comm223\_232\_chi1b (Tensor Branch)}
\author{Code-faithful audit note}
\date{\today}

\begin{document}
\maketitle

This note records the currently implemented equations in
\texttt{src/FactorizedDoubleCommutator\_eths.cc} for the tensor branch inside
\texttt{comm223\_232\_chi1b(Eta,Gamma,Z)}.
It is intentionally code-faithful and may differ from analytic idealization.

\section*{Branch and Conventions}
The documented branch is used when:
\begin{align}
\texttt{tensor\_case} \equiv (\texttt{Eta.GetJRank()}\neq 0)\land(\texttt{Gamma.GetJRank()}=0)\land(\texttt{Z.TwoBody.rank\_T}=0).
\end{align}
Define
\begin{align}
\lambda &= \mathrm{Eta.GetJRank()},
&
\hat\lambda^{-1} &= \frac{1}{\sqrt{2\lambda+1}},
&
\phi(x) &= (-1)^x.
\end{align}

\section*{One-body Intermediates in Tensor Branch}
The code stores
\begin{align}
\mathrm{CHI\_I}(p,q) &\leftarrow \chi^{\epsilon}_{pq},
&
\mathrm{CHI\_II}(p,q) &\leftarrow \chi^{\zeta}_{pq}.
\end{align}

\subsection*{Implemented $\chi^{\epsilon}$ (Omega*Omega, tensor-tensor to scalar)}
For fixed $(p,q)$:
\begin{align}
\chi^{\epsilon}_{pq}
=\frac{1}{2j_p+1}
\sum_{abcJ_0J_1}
\Big(\bar n_a\bar n_b n_c + n_a n_b\bar n_c\Big)
\phi(J_0+J_1+\lambda)
\hat\lambda^{-1}
\Omega_{cpab}^{J_0J_1\lambda}
\Omega_{abcq}^{J_1J_0\lambda},
\end{align}
with loop-domain intersection exactly as in code:
\begin{align}
J_0 &\in
\left[
\max\!\left(\frac{|j_a-j_b|}{2},\frac{|j_c-j_p|}{2}\right),
\min\!\left(\frac{j_a+j_b}{2},\frac{j_c+j_p}{2}\right)
\right],
\\
J_1 &\in
\left[
\max\!\left(\frac{|j_a-j_b|}{2},\frac{|j_c-j_q|}{2}\right),
\min\!\left(\frac{j_a+j_b}{2},\frac{j_c+j_q}{2}\right)
\right],
\end{align}
plus triangle filter $|J_0-J_1|\le\lambda\le J_0+J_1$.

\subsection*{Implemented $\chi^{\zeta}$ (Gamma*Omega, scalar-tensor to tensor)}
For fixed $(p,q)$, code uses
\begin{align}
\chi^{\zeta}_{pq}
=\frac{1}{2}\phi\!\left(\frac{j_q}{2}+\lambda\right)
\sum_{abcJ_0J_1}
\Big(\bar n_a\bar n_b n_c + n_a n_b\bar n_c\Big)
\phi\!\left(J_0+\frac{j_a}{2}\right)
\hat J_0\hat J_1
\sixj{\lambda}{J_1}{J_0}{j_a/2}{j_p/2}{j_q/2}
\Gamma_{apbc}^{J_0J_0 0}
\Omega_{bcaq}^{J_0J_1\lambda},
\end{align}
where $\hat J\equiv\sqrt{2J+1}$,
and the implemented ranges are
\begin{align}
J_0 &\in
\left[
\max\!\left(\frac{|j_b-j_c|}{2},\frac{|j_a-j_p|}{2}\right),
\min\!\left(\frac{j_b+j_c}{2},\frac{j_a+j_p}{2}\right)
\right],
\\
J_1 &\in
\left[
\frac{|j_a-j_q|}{2},\frac{j_a+j_q}{2}
\right],
\end{align}
plus $|J_0-J_1|\le\lambda\le J_0+J_1$.

\section*{Two-body Update: Implemented $\Gamma^{(I)}$ Piece}
The code reuses the original channel-local contraction with
\texttt{GetTBME\_norm(...)} and antisymmetry normalization factors
(\texttt{norm} built from $\sqrt{2}$ and ket/bra phase factors).
In this reused block, all terms multiplying \texttt{CHI\_I}
remain active in tensor case, providing the implemented $\Gamma^{(I)}$-type part.

\section*{Two-body Update: Implemented $\Gamma^{(II)}$ Piece (explicit)}
In tensor case, the old scalar-only \texttt{CHI\_II*Eta} contractions are disabled.
Instead, for each matrix element $(pq;rs)$ in channel $J_0\equiv J$, code adds
\begin{align}
\Gamma^{(II)}_{pqrs} \mathrel{+}= \gammaII_{pqrs},
\end{align}
with
\begin{align}
\gammaII_{pqrs}
&=
\phi\!\left(\frac{j_p}{2}\right)\hat J_0^{-1}
\sum_{aJ_2}
\phi\!\left(J_2+\frac{j_a}{2}\right)
\hat J_2\hat\lambda^{-1}
\sixj{J_0}{J_2}{\lambda}{j_a/2}{j_q/2}{j_p/2}
\chi^{\zeta}_{qa}
\Omega_{pars}^{J_2J_0\lambda}
\\
&\quad-
\phi\!\left(\frac{j_s}{2}\right)\hat J_0^{-1}
\sum_{aJ_2}
\phi\!\left(J_2+\frac{j_a}{2}\right)
\hat J_2\hat\lambda^{-1}
\sixj{J_0}{J_2}{\lambda}{j_a/2}{j_r/2}{j_s/2}
\chi^{\zeta}_{ar}
\Omega_{pqas}^{J_0J_2\lambda},
\end{align}
with triangle filter $|J_0-J_2|\le\lambda\le J_0+J_2$ and per-term
$J_2$ limits from the corresponding pair-coupling intersections used in code:
\begin{align}
J_2:&\ (p,a)\leftrightarrow(r,s) \quad\text{for the first term},
\\
J_2:&\ (a,s)\leftrightarrow(p,q) \quad\text{for the second term}.
\end{align}

\section*{Final Note}
This file mirrors the implemented code path and variable/argument ordering.
It is intended for validation against \texttt{diag2\_compact} term-by-term.

\end{document}
